
Every major LLM pipeline employs data curation, yet optimal filtering parameters remain undetermined. RedPajama, Dolma, and C4 all use perplexity filtering with different thresholds chosen without controlled ablation to justify specific values.

\subsubsection{1.1 The Problem of Curation Calibration}

Current practice treats perplexity thresholds as discrete configuration choices rather than continuous variables with quantifiable dose-response relationships. Pipeline papers compare final configurations while varying multiple factors simultaneously, making it impossible to attribute performance differences to specific curation decisions.

The core methodological gap is the absence of controlled ablation studies that isolate individual curation parameters while holding other factors constant.

\subsubsection{1.2 Key Insight: Dose-Response Relationships}

The central hypothesis is that perplexity filtering exhibits a concave dose-response relationship with benchmark performance. Too permissive filtering includes noise that dilutes gradient signals; too strict filtering removes diversity that enables generalization. This creates an inverted-U response curve with an identifiable optimum.

\subsubsection{1.3 Contributions}

This work makes three contributions:

\begin{enumerate}
\item Demonstration of non-monotonic dose-response relationships between perplexity filtering thresholds and benchmark performance at proof-of-concept scale, with cubic polynomial fit achieving $R^2$ = 0.985.
\end{enumerate}

\begin{enumerate}
\item Validation of the noise-dilution mechanism through convergence dynamics analysis, showing 40\% higher convergence AUC for unfiltered versus moderate filtering.
\end{enumerate}

\begin{enumerate}
\item Scale transfer analysis showing optimal thresholds at 125M transfer to 1B models with ratio 0.85, and CPDR-optimized configuration outperforming RedPajama defaults by 1.32\% at PoC scale.
\end{enumerate}

\textbf{Important Limitation:} All experiments were conducted at proof-of-concept scale (5M--10M tokens) using simulation-based or mock benchmark evaluation. Results are directional; full-scale replication is required before production deployment.
